% §897 takes a kern into the word it is collecting *by subtype*, and anything
% else ends the word outright:
%
%     else if (type(s)=kern_node)and(subtype(s)=normal) then
%       begin hb:=s;
%       ...
%       end
%     else goto done3;
%
% §896's skip loop ahead of the word tests the same way. §1125's `\accent` puts
% a kern of subtype `acc_kern` on each side of its character, so a word with an
% accent in it ends at the accent.
%
% That alone would only shorten the word. What makes it visible is §894's rule
% that there is **one** attempt per glue node — see
% hyphenate-once-per-glue.tex. TeX finds `Gr`, cannot hyphenate two letters, and
% never looks at `unbaum`:
%
%     ref    XX Gr^^?unbaum YY        XX co^^?operation YY
%     ours   XX Gr^^?un-baum YY       XX co^^?-op-er-a-tion YY
%
% `Gr\"unbaum` is the first author cited in 2608.17805v1, and hyphenating it was
% the whole difference between that paper and pdfTeX.
\input plain
\hsize=90pt
\parindent=0pt
\pretolerance=-1
\tolerance=9999
\showboxbreadth=100
\showboxdepth=10
% An accent partway through: the word ends at it, and `Gr` is too short.
\setbox0=\vbox{\noindent XX Gr\"unbaum YY\par}
\showbox0
% The same letters without the accent hyphenate normally.
\setbox0=\vbox{\noindent XX Grunbaum YY\par}
\showbox0
% An accent early in a long word still stops everything behind it.
\setbox0=\vbox{\noindent XX co\"operation YY\par}
\showbox0
% …and the unaccented spelling does hyphenate.
\setbox0=\vbox{\noindent XX cooperation YY\par}
\showbox0
% An accent on the *first* letter leaves nothing before it to hyphenate.
\setbox0=\vbox{\noindent XX \"operation YY\par}
\showbox0
% A normal kern — the kind a font asks for between two letters — does not end
% the word; `\/` is an explicit one and does.
\setbox0=\vbox{\noindent XX opera\/tion YY\par}
\showbox0
\end
